\documentclass[12pt,a4paper]{article}

% --- Кодировки ---
\usepackage[T2A]{fontenc}
\usepackage[utf8]{inputenc}
\usepackage[russian,english]{babel}
\usepackage{cmap}

% --- Геометрия ---
\usepackage[left=2.5cm,right=2cm,top=2cm,bottom=2cm]{geometry}
\usepackage{setspace}
\onehalfspacing

% --- Пакеты ---
\usepackage{graphicx,float,hyperref,enumitem,booktabs,amsmath,amssymb,xcolor,listings,tocloft}
\renewcommand{\contentsname}{}
\renewcommand{\cftsecafterpnum}{\vspace{2pt}}

% --- Компактные списки ---
\setlist{nosep,leftmargin=*}

% --- Листинги ---
\lstset{
    basicstyle=\ttfamily\footnotesize, breaklines=true, frame=single,
    backgroundcolor=\color{gray!5}, rulecolor=\color{gray!50},
    numbers=left, numberstyle=\tiny\color{gray}, numbersep=4pt,
    tabsize=2, showstringspaces=false, aboveskip=4pt, belowskip=4pt,
    language=Python,
}

% --- Гиперссылки ---
\hypersetup{colorlinks=true,linkcolor=blue!50!black,urlcolor=blue!50!black}

% --- Вставка скриншота ---
\newcommand{\img}[4]{%
    \IfFileExists{#1}{%
        \begin{figure}[H]\centering
            \includegraphics[width=#4\textwidth,keepaspectratio]{#1}
            \caption{#2}\label{#3}
        \end{figure}
    }{}%
}

% --- Авторы ---
\newcommand{\AuthorA}{Ясаков Артём}
\newcommand{\AuthorB}{Сарваров Тимур}
\newcommand{\Group}{P3313}
\newcommand{\Teacher}{Жданов Андрей Дмитриевич}
\newcommand{\ImgDir}{/home/yaart/ITMO/3 курс/СИИ/2 модуль/report_images}

% --- Уплотнение ---
\raggedbottom
\setlength{\intextsep}{4pt plus 1pt minus 1pt}
\setlength{\textfloatsep}{6pt plus 1pt minus 1pt}
\setlength{\floatsep}{6pt plus 1pt minus 1pt}
\setlength{\abovecaptionskip}{3pt}
\setlength{\belowcaptionskip}{2pt}

\begin{document}

\begin{titlepage}\centering
{\large Федеральное государственное автономное образовательное\\
учреждение высшего образования\\
«Национальный исследовательский университет ИТМО»\par}
{\large Факультет программной инженерии и компьютерной техники\par}
\vspace{2.5cm}
{\Huge\bfseries Отчёт по Модулю 2\par}
\vspace{0.4cm}
{\Large Классические методы машинного обучения\par}
\vspace{0.4cm}
{\Large Линейная регрессия, \(k\)-NN, деревья решений,\\
логистическая регрессия\par}
\vspace{3.5cm}
\begin{flushright}
\begin{tabular}{rl}
    \textbf{Выполнили:} & \AuthorA \\
                        & \AuthorB \\[6pt]
    \textbf{Группа:}    & \Group \\[12pt]
    \textbf{Проверил:}  & \Teacher \\
\end{tabular}
\end{flushright}
\vfill
{\large Санкт-Петербург, 2026 г.\par}
\end{titlepage}

\tableofcontents
\newpage

% ============================================================
\section*{Аннотация}
\addcontentsline{toc}{section}{Аннотация}

В рамках Модуля~2 дисциплины «Системы искусственного интеллекта»
реализованы «с нуля» четыре классических алгоритма машинного обучения:
линейная регрессия (California Housing), метод \(k\)-ближайших соседей
(Wine), дерево решений ID3 (Mushroom) и логистическая регрессия
(Titanic). Все реализации выполнены на Python с использованием только
NumPy и Pandas (кроме дерева решений, использующего исключительно
стандартную библиотеку). Проведён сравнительный анализ методов по
метрикам качества, исследовано влияние гиперпараметров и наборов
признаков. Сформулированы практические рекомендации по выбору метода
в зависимости от типа задачи, размера данных и требований
к интерпретируемости.

\medskip
\noindent\textbf{Ключевые слова:} машинное обучение, линейная регрессия,
\(k\)-ближайших соседей, дерево решений, логистическая регрессия,
градиентный спуск, метод Ньютона, классификация, регрессия,
гиперпараметры.

\newpage

% ============================================================
\section{Введение}

\subsection{Цели модуля}

Модуль~2 направлен на освоение фундаментальных алгоритмов машинного
обучения через их самостоятельную реализацию. В ходе модуля решаются
следующие задачи:

\begin{itemize}
    \item Линейная регрессия — предсказание непрерывной величины
          (стоимость жилья) по набору числовых и категориальных
          признаков.
    \item \(k\)-ближайших соседей — классификация объектов (сорт вина)
          без явного обучения модели.
    \item Дерево решений — интерпретируемая классификация (съедобность
          грибов) с автоматическим построением правил.
    \item Логистическая регрессия — бинарная классификация (выживаемость
          пассажиров Титаника) с вероятностной интерпретацией.
\end{itemize}

\subsection{Критерии выбора метода}

Выбор алгоритма определяется несколькими ключевыми факторами:

\begin{itemize}
    \item \textbf{Тип целевой переменной.} Непрерывная — регрессия;
          категориальная — классификация.
    \item \textbf{Размер и качество данных.} \(k\)-NN чувствителен
          к размерности и требует масштабирования; деревья решений
          устойчивы к пропускам и выбросам.
    \item \textbf{Интерпретируемость.} Линейные модели и деревья
          позволяют понять вклад каждого признака; \(k\)-NN не даёт
          такой возможности.
    \item \textbf{Стоимость ошибок.} В задачах классификации выбор
          между Precision и Recall зависит от цены ложноположительных
          и ложноотрицательных срабатываний.
\end{itemize}

\subsection{Структура отчёта}

Отчёт последовательно описывает четыре лабораторные работы, каждая
из которых содержит: постановку задачи, описание данных, принцип
работы метода с псевдокодом, результаты экспериментов и рекомендации
по применению. Завершается отчёт сравнительным анализом методов
и общими выводами.

\newpage

% ============================================================
\section{Данные и общая методология}

\subsection{Выбранные датасеты}

Для каждой лабораторной работы использован отдельный датасет,
подобранный под специфику метода (таблица~\ref{tab:datasets}).

\begin{table}[H]\centering\small
\caption{Обзор датасетов}\label{tab:datasets}
\begin{tabular}{llp{3.5cm}ccc}\toprule
Лаб. & Метод & Датасет & Объектов & Признаков & Классов \\ \midrule
3 & Линейная регрессия & California Housing & 20\,640 & 8 + 1 кат. & — \\
4 & \(k\)-NN & Wine & 178 & 13 & 3 \\
5 & Дерево решений & Mushroom & 8\,124 & 22 кат. & 2 \\
6 & Логист. регрессия & Titanic & 891 & 7 + 2 кат. & 2 \\ \bottomrule
\end{tabular}
\end{table}

\subsection{Предобработка данных}

Во всех работах применялась единая методология:

\begin{enumerate}
    \item \textbf{Пропуски:} заполнение медианой (числовые) или
          исключение признака (категориальные с большим числом
          пропусков, например, \texttt{deck} в Titanic — 688 из 891).
    \item \textbf{Категориальные признаки:} one-hot кодирование
          с~\texttt{drop\_first} для линейных моделей; для дерева
          решений используются исходные категории.
    \item \textbf{Масштабирование:} Z-нормировка по статистикам
          обучающей выборки для методов, чувствительных к масштабу
          (\(k\)-NN, логистическая регрессия).
    \item \textbf{Разбиение:} train/test 80/20 со стратификацией
          (для классификации) или случайное (для регрессии).
    \item \textbf{Random seed:} 42 — фиксирован для воспроизводимости.
\end{enumerate}

\subsection{Инфраструктура воспроизводимости}

\begin{itemize}
    \item Язык: Python 3.13.
    \item Зависимости: NumPy, Pandas, Matplotlib, Seaborn, scikit-learn
          (только для \texttt{train\_test\_split} и
          \texttt{StandardScaler} в~лабах 4 и~6; дерево решений —
          исключительно стандартная библиотека).
    \item Запуск: \texttt{python3 lab\_*.py} из соответствующей
          директории.
\end{itemize}

\newpage

% ============================================================
\section{Лабораторная работа 3. Линейная регрессия}

\subsection{Постановка задачи}

Предсказать медианную стоимость жилья (\texttt{median\_house\_value})
в~районах Калифорнии по набору социально-экономических
и~географических признаков.

\subsection{Описание данных}

\begin{itemize}
    \item \textbf{Источник:} California Housing Prices (A.~Géron,
          Hands-On Machine Learning).
    \item \textbf{Объектов:} 20\,640 районов.
    \item \textbf{Признаки:} 8 числовых (\texttt{longitude},
          \texttt{latitude}, \texttt{housing\_median\_age},
          \texttt{total\_rooms}, \texttt{total\_bedrooms},
          \texttt{population}, \texttt{households},
          \texttt{median\_income}) + 1 категориальный
          (\texttt{ocean\_proximity}, 5 значений).
    \item \textbf{Целевая:} \texttt{median\_house\_value} (доллары США).
    \item \textbf{Пропуски:} 207 в \texttt{total\_bedrooms} (заполнены
          медианой 435.0).
\end{itemize}

\img{\ImgDir/lr_target_hist.png}{Распределение целевой переменной}{fig:lr-hist}{0.85}
\img{\ImgDir/lr_income_vs_target.png}{Доход vs Стоимость жилья}{fig:lr-scatter}{0.85}

\subsection{Зачем метод}

Линейная регрессия — базовый метод для задач с непрерывной целевой
переменной. Модель предполагает линейную зависимость:
\(\hat{y} = \beta_0 + \sum_{j=1}^{d} \beta_j x_j\).

\textbf{Условия применимости:} линейная связь признаков с целевой,
отсутствие сильной мультиколлинеарности, нормальность остатков.
\textbf{Ограничения:} не улавливает нелинейные зависимости,
чувствительна к выбросам.

\subsection{Принцип работы}

Метод наименьших квадратов (МНК) минимизирует сумму квадратов
остатков: \(\hat{\beta} = \arg\min_{\beta} \|y - X\beta\|_2^2\).
Аналитическое решение через нормальные уравнения:
\(\hat{\beta} = (X^T X)^{-1} X^T y\).

\fbox{\begin{minipage}{0.95\textwidth}
\textbf{Псевдокод:} Линейная регрессия (МНК)
\begin{enumerate}
    \item \textbf{Вход:} \(X \in \mathbb{R}^{n \times d}\), \(y \in \mathbb{R}^n\).
    \item Добавить столбец единиц: \(X_{\text{aug}} = [\mathbf{1} \mid X]\).
    \item Составить нормальные уравнения: \(A = X_{\text{aug}}^T X_{\text{aug}}\),
          \(b = X_{\text{aug}}^T y\).
    \item Решить \(A\beta = b\) методом Гаусса с выбором главного элемента.
    \item \textbf{Выход:} \(\beta = [\beta_0, \beta_1, \dots, \beta_d]\).
\end{enumerate}
\end{minipage}}

\medskip
\textbf{Ключевой фрагмент — решение СЛАУ методом Гаусса (без \texttt{np.linalg}):}

\begin{lstlisting}[caption={Метод Гаусса с выбором главного элемента},label={lst:gauss}]
def gaussian_solve(A, b):
    A, b = np.array(A, dtype=float, copy=True), np.array(b, dtype=float, copy=True)
    n = len(b)
    for k in range(n):
        pivot = k + int(np.argmax(np.abs(A[k:, k])))
        if pivot != k: A[[k, pivot]] = A[[pivot, k]]; b[[k, pivot]] = b[[pivot, k]]
        A[k, k:] /= A[k, k]; b[k] /= A[k, k]
        for i in range(n):
            if i == k: continue
            factor = A[i, k]; A[i, k:] -= factor * A[k, k:]; b[i] -= factor * b[k]
    return b
\end{lstlisting}

\subsection{Эксперименты}

Построены 4 модели с последовательным расширением набора признаков:

\begin{enumerate}
    \item \textbf{Модель 1} — только \texttt{median\_income}.
    \item \textbf{Модель 2} — \texttt{median\_income},
          \texttt{housing\_median\_age}, \texttt{latitude},
          \texttt{longitude}.
    \item \textbf{Модель 3} — все 8 числовых признаков + one-hot
          кодирование \texttt{ocean\_proximity} (4 бинарных признака,
          базовая категория: \texttt{<1H OCEAN}).
    \item \textbf{Бонус} — модель 3 + синтетический признак
          \texttt{bedrooms\_per\_room}.
\end{enumerate}

\begin{table}[H]\centering
\caption{Результаты линейной регрессии}\label{tab:lr-results}
\begin{tabular}{lccc}\toprule
Модель & Признаков & R² (train) & R² (test) \\ \midrule
Модель 1 (только доход)    & 1  & 0.4748 & 0.4683 \\
Модель 2 (+гео+возраст)    & 4  & 0.5954 & 0.5886 \\
Модель 3 (все признаки)    & 12 & 0.6462 & 0.6423 \\
Бонус (+синтетический)     & 13 & 0.6487 & \textbf{0.6446} \\ \bottomrule
\end{tabular}
\end{table}

\img{\ImgDir/lr_r2_comparison.png}{Сравнение R² моделей}{fig:lr-r2}{0.85}
\img{\ImgDir/lr_correlation.png}{Корреляционная матрица числовых признаков}{fig:lr-corr}{0.9}

\textbf{Анализ:}
\begin{itemize}
    \item \texttt{median\_income} — доминирующий признак: одна лишь
          эта переменная даёт \(R^2 \approx 0.47\).
    \item Добавление географических координат и возраста жилья
          увеличивает \(R^2\) на \(\sim 0.12\).
    \item Полный набор с категориальными признаками даёт \(R^2 \approx
          0.64\); синтетический признак даёт минимальный прирост
          (\(+0.002\)).
    \item Переобучение отсутствует: разница train/test \(R^2\)
          минимальна во всех моделях.
\end{itemize}

\subsection{Когда использовать}

\begin{itemize}
    \item Прогнозирование цен недвижимости, аренды.
    \item Оценка влияния факторов на экономические показатели.
    \item Быстрый baseline для задач регрессии перед применением
          сложных моделей.
    \item Задачи, где требуется интерпретация коэффициентов (на сколько
          изменится целевая при изменении признака на единицу).
\end{itemize}

\newpage

% ============================================================
\section{Лабораторная работа 4. Метод \(k\)-ближайших соседей}

\subsection{Постановка задачи}

Классифицировать сорт вина по 13 химическим показателям. Многоклассовая
задача (3 сорта винограда).

\subsection{Описание данных}

\begin{itemize}
    \item \textbf{Источник:} \texttt{sklearn.datasets.load\_wine} (UCI
          Wine).
    \item \textbf{Объектов:} 178 образцов (сбалансированы: 59, 71, 48).
    \item \textbf{Признаки:} 13 числовых (alcohol, malic\_acid, ash,
          alcalinity\_of\_ash, magnesium, total\_phenols, flavanoids,
          nonflavanoid\_phenols, proanthocyanins, color\_intensity,
          hue, od280/od315, proline).
    \item \textbf{Пропуски:} отсутствуют.
\end{itemize}

\img{\ImgDir/knn_3d_scatter.png}{3D-визуализация: alcohol, flavanoids, color\_intensity}{fig:knn-3d}{0.8}

\subsection{Зачем метод}

\(k\)-NN — метод «ленивого обучения»: модель не имеет явной фазы
обучения, а запоминает все обучающие примеры. Классификация
выполняется по принципу «похожие объекты принадлежат одному классу».

\textbf{Требования:} обязательное масштабирование признаков
(чувствителен к масштабу), неэффективен при высокой размерности
(«проклятие размерности»). \textbf{Ограничения:} медленный на этапе
предсказания (\(O(n_{\text{train}} \cdot d)\) на один объект),
чувствителен к шуму и неинформативным признакам.

\subsection{Принцип работы}

\fbox{\begin{minipage}{0.95\textwidth}
\textbf{Псевдокод:} \(k\)-ближайших соседей
\begin{enumerate}
    \item \textbf{Вход:} обучающая выборка \((X_{\text{train}}, y_{\text{train}})\),
          тестовый объект \(x\), число соседей \(k\).
    \item Вычислить евклидовы расстояния до всех обучающих объектов:
          \(d_i = \|x - x_i\|_2\).
    \item Найти \(k\) объектов с наименьшими \(d_i\).
    \item Класс — мажоритарное голосование:
          \(\hat{y} = \arg\max_c \sum_{i \in \text{top-}k} \mathbb{1}[y_i = c]\).
    \item \textbf{Выход:} предсказанный класс \(\hat{y}\).
\end{enumerate}
\end{minipage}}

\medskip
Для ускорения расстояния вычисляются векторизованно через тождество:
\(\|a - b\|^2 = \|a\|^2 + \|b\|^2 - 2 a \cdot b\), что позволяет
вычислить всю матрицу расстояний \((n_{\text{test}} \times
n_{\text{train}})\) за одно матричное умножение.

\begin{lstlisting}[caption={Векторизованное вычисление расстояний и голосование},label={lst:knn}]
# Матрица расстояний (n_test x n_train) без циклов
dists = (np.sum(X_test**2, axis=1)[:, None]
       + np.sum(X_train**2, axis=1)[None, :]
       - 2 * np.dot(X_test, X_train.T))
# Для каждого тестового объекта
top_k = y_train[np.argsort(dists[i])[:k]]
pred = np.argmax(np.bincount(top_k, minlength=n_classes))
\end{lstlisting}

\subsection{Эксперименты}

Построены две модели с разными наборами признаков:
\begin{enumerate}
    \item \textbf{Модель 1 (случайные):} 6 случайно отобранных
          признаков (\texttt{od280/od315}, \texttt{color\_intensity},
          \texttt{alcohol}, \texttt{proanthocyanins},
          \texttt{total\_phenols}, \texttt{ash}).
    \item \textbf{Модель 2 (фиксированные):} 5 наиболее информативных
          признаков (\texttt{alcohol}, \texttt{flavanoids},
          \texttt{color\_intensity}, \texttt{proline},
          \texttt{total\_phenols}).
\end{enumerate}

Исследованы \(k \in \{1, 3, 5, 7, 10, 15\}\).

\begin{table}[H]\centering
\caption{Точность \(k\)-NN при разных \(k\)}
\begin{tabular}{ccc}\toprule
\(k\) & Модель 1 (случайные) & Модель 2 (фиксированные) \\ \midrule
1  & 0.8889 & 0.9630 \\
3  & 0.8889 & 0.9630 \\
5  & 0.9074 & 0.9630 \\
7  & \textbf{0.9259} & 0.9630 \\
10 & 0.8889 & \textbf{0.9815} \\
15 & 0.8889 & 0.9444 \\ \bottomrule
\end{tabular}
\end{table}

\img{\ImgDir/knn_acc_vs_k.png}{Зависимость точности от \(k\)}{fig:knn-acc}{0.85}
\img{\ImgDir/knn_confusion_matrices.png}{Матрицы ошибок (верхний ряд — Модель 1, нижний — Модель 2)}{fig:knn-cm}{1.0}

\textbf{Анализ:}
\begin{itemize}
    \item Модель 2 стабильно превосходит Модель 1 за счёт включения
          \texttt{flavanoids} и \texttt{proline} — наиболее
          разделяющих признаков.
    \item Оптимальное \(k = 10\) для Модели 2 (Accuracy = 0.9815).
    \item При \(k = 15\) точность снижается из-за «размытия» границ
          решающей области.
    \item Класс 0 распознаётся практически безошибочно; ошибки
          сосредоточены между классами 1 и 2 (матрицы ошибок
          подтверждают).
\end{itemize}

\subsection{Когда использовать}

\begin{itemize}
    \item Быстрый baseline без настройки (нет фазы обучения).
    \item Многоклассовая классификация с небольшим числом признаков.
    \item Рекомендательные системы (поиск похожих товаров/фильмов).
    \item Обнаружение аномалий (по расстоянию до соседей).
\end{itemize}

\newpage

% ============================================================
\section{Лабораторная работа 5. Деревья решений}

\subsection{Постановка задачи}

Классифицировать грибы на съедобные (\texttt{e}) и ядовитые
(\texttt{p}) по набору категориальных признаков. Задача бинарной
классификации с высокими требованиями к интерпретируемости
(цена ошибки — жизнь человека).

\subsection{Описание данных}

\begin{itemize}
    \item \textbf{Источник:} UCI Mushroom (Agaricus-Lepiota).
    \item \textbf{Объектов:} 8\,124 образца (4\,208 съедобных,
          3\,916 ядовитых).
    \item \textbf{Признаки:} 22 категориальных (cap-shape, cap-color,
          odor, gill-size, spore-print-color, population, habitat
          и~др.).
    \item \textbf{Пропуски:} 2\,480 в \texttt{stalk-root}
          (обрабатываются как отдельная категория \texttt{?}).
\end{itemize}

\img{\ImgDir/dt_class_dist.png}{Распределение классов}{fig:dt-dist}{0.6}

\subsection{Зачем метод}

Деревья решений привлекательны своей \textbf{интерпретируемостью}:
модель представляет собой набор правил «если… то…», понятных
человеку. Метод автоматически выявляет нелинейные взаимодействия,
не требует масштабирования и устойчив к выбросам.

\textbf{Ограничения:} склонность к переобучению (глубокие деревья),
чувствительность к малым изменениям данных, жадная стратегия
построения не гарантирует глобального оптимума.

\subsection{Принцип работы}

Алгоритм ID3 строит дерево рекурсивно, выбирая на каждом шаге признак,
дающий максимальный прирост информации:
\[
IG(S, A) = H(S) - \sum_{v \in \text{values}(A)} \frac{|S_v|}{|S|} H(S_v),
\]
где \(H(S) = -\sum_{c} p_c \log_2 p_c\) — энтропия Шеннона.

\fbox{\begin{minipage}{0.95\textwidth}
\textbf{Псевдокод:} Построение дерева (ID3)
\begin{enumerate}
    \item \textbf{Вход:} выборка \(S\), список признаков \(F\).
    \item Если все объекты одного класса или \(F\) пуст — создать
          лист с вероятностью класса \(P(\text{positive})\).
    \item Для каждого признака \(A \in F\) вычислить
          \(\text{IG}(S, A)\).
    \item Выбрать \(A^* = \arg\max_A \text{IG}(S, A)\).
    \item Для каждого значения \(v \in A^*\) рекурсивно построить
          поддерево на \(S_v\) с признаками \(F \setminus \{A^*\}\).
    \item \textbf{Выход:} узел с признаком \(A^*\) и дочерними
          поддеревьями.
\end{enumerate}
\end{minipage}}

\medskip
Из 22 признаков случайно отбираются \(\lfloor\sqrt{22}\rfloor = 4\)
(аналог Random Forest). Библиотеки: только \texttt{csv}, \texttt{math},
\texttt{random} — никаких NumPy/Pandas/sklearn.

\begin{lstlisting}[caption={Построение дерева ID3 — выбор лучшего признака},label={lst:tree}]
def build_tree(rows, features):
    prob = sum(r['class'] == 'p' for r in rows) / len(rows)
    if prob in (0, 1) or not features:
        return Node(len(rows), prob)
    best_gain, best_feature = -math.inf, None
    for f in features:
        groups = {}; [groups.setdefault(r[f], []).append(r) for r in rows]
        gain = entropy(rows) - sum(
            len(g)/len(rows)*entropy(g) for g in groups.values())
        if gain > best_gain + 1e-12:
            best_gain, best_feature, best_groups = gain, f, groups
    node = Node(len(rows), prob); node.feature = best_feature
    node.children = {v: build_tree(g, [x for x in features if x != best_feature])
                     for v, g in sorted(best_groups.items())}
    return node
\end{lstlisting}

\subsection{Эксперименты}

Отобранные признаки: \texttt{population}, \texttt{bruises},
\texttt{cap-shape}, \texttt{gill-color}. Корень дерева —
\texttt{gill-color} (цвет пластинок гриба).

\begin{table}[H]\centering\small
\caption{Характеристики и метрики дерева}\label{tab:dt}
\begin{tabular}{lclc}\toprule
Параметр & Значение & Метрика (test) & Значение \\ \midrule
Узлов / листьев & 141 / 96 & Accuracy & 0.9335 \\
Глубина & 4 & Precision & 0.9592 \\
Корень & gill-color & Recall & 0.9004 \\
Макс. потомков & 12 & ROC-AUC & 0.9828 \\ \bottomrule
\end{tabular}
\end{table}

\img{\ImgDir/dt_tree_viz.png}{Визуализация построенного дерева (первые уровни)}{fig:dt-tree}{0.95}
\img{\ImgDir/dt_metrics.png}{Сравнение метрик train/test}{fig:dt-metrics}{0.8}
\img{\ImgDir/dt_roc.png}{ROC-кривая на тестовой выборке}{fig:dt-roc}{0.7}

\textbf{Анализ:}
\begin{itemize}
    \item Дерево глубиной 4 с 141 узлом — компактная и
          интерпретируемая модель. Ограничение глубины выступает
          как неявная регуляризация (прунинг).
    \item Наиболее информативный признак — \texttt{gill-color}
          (цвет пластинок гриба).
    \item Precision 0.96, Recall 0.90 — модель осторожна:
          предпочитает пропустить съедобный гриб, чем ошибочно назвать
          ядовитый съедобным (что соответствует цене ошибки).
    \item ROC-AUC = 0.983 — отличная разделяющая способность.
    \item Разница train/test минимальна — переобучения нет.
\end{itemize}

\subsection{Когда использовать}

\begin{itemize}
    \item Медицинская диагностика (понятные правила).
    \item Кредитный скоринг (требуется объяснение решения).
    \item Задачи с категориальными признаками и пропусками.
    \item Предварительный анализ данных (feature importance).
\end{itemize}

\newpage

% ============================================================
\section{Лабораторная работа 6. Логистическая регрессия}

\subsection{Постановка задачи}

Предсказать, выжил ли пассажир Титаника (бинарная классификация),
по~набору признаков: класс билета, пол, возраст, количество
родственников, стоимость билета, порт посадки.

\subsection{Описание данных}

\begin{itemize}
    \item \textbf{Источник:} \texttt{seaborn.load\_dataset('titanic')}.
    \item \textbf{Объектов:} 891 пассажир (549 погибших / 342 выживших;
          классы не сбалансированы — 61.6\% / 38.4\%).
    \item \textbf{Признаки после обработки:} 9 числовых: sex, age,
          sibsp, parch, fare, embarked\_Q/S, pclass\_2/3.
    \item \textbf{Пропуски:} age (177), embarked (2) — заполнены
          медианой.
\end{itemize}

\img{\ImgDir/lr_survival_factors.png}{Факторы выживаемости}{fig:lr-factors}{0.85}
\img{\ImgDir/feature_distributions.png}{Распределение признаков датасета Titanic}{fig:lr-dist}{0.9}

\subsection{Зачем метод}

Логистическая регрессия — основной метод бинарной классификации,
дающий вероятностную интерпретацию. В отличие от линейной регрессии,
предсказывает вероятность принадлежности к классу через сигмоиду:
\[
P(y = 1 \mid x) = \sigma(w^T x + b) = \frac{1}{1 + e^{-(w^T x + b)}}.
\]

\textbf{Интерпретация коэффициентов:} \(\beta_j\) показывает изменение
логарифма шансов (\(\log\frac{P}{1-P}\)) при изменении \(x_j\) на
единицу. \textbf{Ограничения:} предполагает линейную разделимость
в~логит-пространстве, чувствительна к масштабу признаков.

\subsection{Принцип работы}

Обучение сводится к минимизации log loss:
\[
\mathcal{L}(w, b) = -\frac{1}{n} \sum_{i=1}^{n}
\bigl[ y_i \ln \hat{y}_i + (1 - y_i) \ln (1 - \hat{y}_i) \bigr],
\]
где \(\hat{y}_i = \sigma(w^T x_i + b)\).

\fbox{\begin{minipage}{0.95\textwidth}
\textbf{Псевдокод:} Логистическая регрессия (градиентный спуск)
\begin{enumerate}
    \item \textbf{Вход:} \(X_{\text{train}}\), \(y_{\text{train}}\), \(\eta\), \(N\).
    \item Инициализировать \(w = 0\), \(b = 0\).
    \item \textbf{Повторить} \(N\) раз:
          \(\hat{y} = \sigma(Xw + b)\);
          \(\nabla_w = \frac{1}{n} X^T (\hat{y} - y)\);
          \(\nabla_b = \frac{1}{n} \sum (\hat{y}_i - y_i)\);
          \(w := w - \eta \nabla_w\); \(b := b - \eta \nabla_b\).
    \item \textbf{Выход:} \(w, b\).
\end{enumerate}
\end{minipage}}

\medskip
Реализован также \textbf{метод Ньютона}: \(\theta := \theta - H^{-1} \nabla \mathcal{L}(\theta)\),
где \(H = \frac{1}{n} X_{\text{aug}}^T S X_{\text{aug}} + \lambda I\) —
регуляризованная матрица Гессе, \(S = \operatorname{diag}(\hat{y}_i(1-\hat{y}_i))\).

\begin{lstlisting}[caption={Градиентный спуск и метод Ньютона},label={lst:lr}]
# Градиентный спуск
class LogisticRegressionGD:
    def fit(self, X, y):
        self.weights = np.zeros((n_features, 1)); self.bias = 0.0
        for _ in range(self.n_iter):
            y_pred = self._sigmoid(np.dot(X, self.weights) + self.bias)
            self.weights -= self.lr * (1/n) * np.dot(X.T, (y_pred - y))
            self.bias    -= self.lr * (1/n) * np.sum(y_pred - y)
        return self

# Метод Ньютона
grad = np.vstack([grad_w, [[grad_b]]])
S = np.diag((y_pred * (1 - y_pred)).flatten())
X_aug = np.hstack([X, np.ones((n, 1))])
H = (1/n) * X_aug.T @ S @ X_aug + 1e-5 * np.eye(n_features + 1)
delta = np.linalg.solve(H, grad)
self.weights -= delta[:-1].reshape(-1, 1); self.bias -= delta[-1, 0]
\end{lstlisting}

\subsection{Эксперименты}

Исследованы гиперпараметры: learning rate
\(\eta \in \{0.001, 0.01, 0.05, 0.1, 0.5, 1.0\}\),
количество итераций \(N \in \{100, 500, 1000, 5000, 10000\}\),
метод оптимизации (градиентный спуск и метод Ньютона).

\begin{table}[H]\centering\small
\caption{Влияние learning rate (GD, \(N = 5000\))}\label{tab:lr-lr}
\begin{tabular}{ccccc}\toprule
\(\eta\) & Accuracy & Precision & Recall & F1 \\ \midrule
0.001 & 0.7765 & 0.7302 & 0.6667 & 0.6970 \\
0.01  & 0.7989 & 0.7797 & 0.6667 & 0.7188 \\
0.05  & 0.8045 & 0.7931 & 0.6667 & 0.7244 \\
\textbf{0.1} & \textbf{0.8045} & \textbf{0.7931} & 0.6667 & \textbf{0.7244} \\
0.5   & 0.8045 & 0.7931 & 0.6667 & 0.7244 \\
1.0   & 0.8045 & 0.7931 & 0.6667 & 0.7244 \\ \bottomrule
\end{tabular}
\end{table}

\begin{table}[H]\centering\small
\caption{Влияние количества итераций (GD, \(\eta = 0.1\))}\label{tab:lr-niter}
\begin{tabular}{ccccc}\toprule
\(N\) & Accuracy & Precision & Recall & F1 \\ \midrule
100   & 0.7989 & 0.7797 & 0.6667 & 0.7188 \\
500   & 0.7989 & 0.7797 & 0.6667 & 0.7188 \\
1000  & 0.8045 & 0.7931 & 0.6667 & 0.7244 \\
5000  & 0.8045 & 0.7931 & 0.6667 & 0.7244 \\
10000 & \textbf{0.8045} & \textbf{0.7931} & 0.6667 & \textbf{0.7244} \\ \bottomrule
\end{tabular}
\end{table}

\begin{table}[H]\centering\small
\caption{Метод Ньютона — сходимость за 3 итерации}
\begin{tabular}{cccc}\toprule
Итераций & Accuracy & Precision & Recall / F1 \\ \midrule
3 и более & 0.8045 & 0.7931 & 0.6667 / 0.7244 \\ \bottomrule
\end{tabular}
\end{table}

\img{\ImgDir/hyperparameter_study.png}{Зависимость метрик от гиперпараметров}{fig:lr-hyper}{0.9}

\img{\ImgDir/heatmap_summary.png}{Тепловая карта Accuracy (GD) и сводная таблица}{fig:lr-heat}{0.9}

\textbf{Анализ:}
\begin{itemize}
    \item При \(\eta = 0.001\) модель сходится медленно
          (Accuracy \(\approx 0.78\)). Оптимальный диапазон:
          \(\eta \in [0.01, 0.1]\).
    \item При \(\eta \geq 0.5\) расходимости не происходит благодаря
          нормализации, но на более сложных датасетах возможна
          нестабильность.
    \item После \(N = 1000\) итераций наступает плато; дальнейшее
          обучение не улучшает метрики.
    \item Метод Ньютона сходится за 3--5 итераций против \(\sim 1000\)
          для GD, но требует \(O(d^3)\) на обращение Гессе.
    \item Precision (0.79) выше Recall (0.67) — модель осторожна
          с~предсказанием класса «выжил», что объясняется дисбалансом
          классов (61.6\% погибших).
\end{itemize}

\subsection{Финальная модель}

С оптимальными параметрами (\(\eta = 0.1\), \(N = 10000\), GD):

\begin{table}[H]\centering
\caption{Финальные метрики}\label{tab:lr-final}
\begin{tabular}{lclc}\toprule
Метрика & Значение & Элемент & Значение \\ \midrule
Accuracy  & 0.8045 & TN & 98 \\
Precision & 0.7931 & FP & 12 \\
Recall    & 0.6667 & FN & 23 \\
F1-Score  & 0.7244 & TP & 46 \\ \bottomrule
\end{tabular}
\end{table}

\img{\ImgDir/final_model.png}{Кривая обучения и матрица ошибок}{fig:lr-final}{0.9}

Модель правильно классифицирует 98 из 110 погибших (Recall \(\approx
0.89\)), но лишь 46 из 69 выживших (Recall \(\approx 0.67\)). Это
консервативное поведение соответствует априорному распределению
классов.

\subsection{Когда использовать}

\begin{itemize}
    \item Бинарная классификация с требованием вероятностного выхода.
    \item Задачи, где важна интерпретация коэффициентов (odds ratios).
    \item Медицинские исследования (вероятность заболевания).
    \item Базовый классификатор перед переходом к сложным моделям.
\end{itemize}

\newpage

% ============================================================
\section{Сравнение методов}

\subsection{Сводная таблица}

\begin{table}[H]\centering\small
\caption{Сравнительная характеристика методов}\label{tab:compare}
\begin{tabular}{p{2.5cm}llll}\toprule
Критерий & Лин. регрессия & \(k\)-NN & Дерево решений & Логист. регрессия \\ \midrule
Тип задачи & Регрессия & Классификация & Классификация & Классификация \\
Обучение & Аналитическое & Ленивое & Рекурсивное & Итеративное \\
Масштабирование & Желательно & Обязательно & Не требуется & Обязательно \\
Интерпретируемость & Высокая & Низкая & Очень высокая & Высокая \\
Чувств. к выбросам & Высокая & Средняя & Низкая & Средняя \\
Чувств. к размерности & Низкая & Высокая & Средняя & Низкая \\
Нелинейность & Нет & Да & Да & Нет \\
Пропуски & Мешают & Мешают & Обрабатывает & Мешают \\
Скорость predict & \(O(d)\) & \(O(nd)\) & \(O(\text{depth})\) & \(O(d)\) \\
Вероятностный выход & Нет & Нет & Да & Да \\ \bottomrule
\end{tabular}
\end{table}

\subsection{Рекомендации «если…, то…»}

\begin{enumerate}
    \item \textbf{Мало признаков + нужна интерпретация}
          \(\rightarrow\) линейная или логистическая регрессия.
    \item \textbf{Сложные нелинейные зависимости}
          \(\rightarrow\) дерево решений.
    \item \textbf{Нет времени на обучение / быстрый baseline}
          \(\rightarrow\) \(k\)-NN.
    \item \textbf{Категориальные признаки / пропуски в данных}
          \(\rightarrow\) дерево решений.
    \item \textbf{Вероятностный прогноз для бинарного исхода}
          \(\rightarrow\) логистическая регрессия.
    \item \textbf{Высокая размерность (\(d > 100\))}
          \(\rightarrow\) линейная/логистическая регрессия.
    \item \textbf{Цена ошибки асимметрична} \(\rightarrow\) метод
          с~настраиваемым порогом + оценка по Precision/Recall.
\end{enumerate}

\subsection{Итоговые метрики}

\begin{table}[H]\centering
\caption{Лучшие результаты модуля}\label{tab:final-all}
\begin{tabular}{llc}\toprule
Метод & Метрика & Значение \\ \midrule
Линейная регрессия  & R² & 0.6446 \\
\(k\)-NN (Wine)     & Accuracy & 0.9815 \\
Дерево решений      & ROC-AUC & 0.9828 \\
Логист. регрессия   & F1-Score & 0.7244 \\ \bottomrule
\end{tabular}
\end{table}

\newpage

% ============================================================
\section{Заключение}

В ходе выполнения Модуля~2 освоены четыре фундаментальных алгоритма
машинного обучения, каждый из которых реализован «с нуля» без опоры
на высокоуровневые ML-фреймворки. Основные выводы:

\begin{enumerate}
    \item \textbf{Линейная регрессия} — простой и интерпретируемый
          метод для задач регрессии. Ключевой фактор качества —
          выбор информативных признаков (\texttt{median\_income} даёт
          почти половину объяснённой дисперсии).
    \item \textbf{\(k\)-NN} — эффективен при малой размерности и
          хорошем масштабировании. Качество критически зависит от
          набора признаков: переход от случайных к информативным
          повысил Accuracy с~0.93 до~0.98.
    \item \textbf{Дерево решений} — лучший выбор, когда требуется
          интерпретируемость. На датасете Mushroom модель глубиной~4
          достигла ROC-AUC 0.983, при этом правила легко читаются
          человеком.
    \item \textbf{Логистическая регрессия} — надёжный baseline для
          бинарной классификации с вероятностным выходом. Метод
          Ньютона сходится на порядки быстрее градиентного спуска,
          но практический выигрыш проявляется только при больших~\(d\).
    \item \textbf{Выбор метода} определяется типом задачи, размером
          и~качеством данных, требованиями к интерпретируемости и
          допустимой ценой ошибок. Не существует универсального
          алгоритма — необходимо сопоставлять свойства метода
          с~особенностями конкретной задачи.
\end{enumerate}

\newpage

\section*{Использование искусственного интеллекта}
\addcontentsline{toc}{section}{Использование искусственного интеллекта}

В ходе работы использовались инструменты ИИ (Claude, ChatGPT) для
помощи в написании кода, отладке алгоритмов и оформлении отчёта.
Конкретно: реализация метода Ньютона и ID3 с нуля, векторизация
расстояний в k-NN, подбор формул для отчёта, генерация графиков
и~вёрстка \LaTeX. Все решения осмыслены и проверены авторами.

\end{document}